\chapter{Arithmetische Funktionen}
\label{b10guide:ch:arithmetic-functions}

Dieses Kapitel behandelt drei konkrete Rechengrößen positiver
natürlicher Zahlen: Teileranzahl, Teilersumme und Eulers
$\varphi$-Funktion. Ihre Definitionen verwenden nur endliche
natürliche Summen und Fallunterscheidungen.

\section{Teileranzahl und Teilersumme}

\begin{definition}[Teileranzahl und Teilersumme]
\label{b10guide:def:divisor-functions}
Für $n,d\in\mathbb N_+$ sei
\[
 \delta_n(d)=
 \begin{cases}
  1,&d\mid n,\\
  0,&d\nmid n.
 \end{cases}
\]
Die \emph{Teileranzahl} $\tau(n)$ und die \emph{Teilersumme}
$\sigma(n)$ sind durch
\[
 \tau(n)=\sum_{i<n}\delta_n(i+1),\qquad
 \sigma(n)=\sum_{i<n}(i+1)\delta_n(i+1)
\]
festgelegt. Damit erfasst $\tau(n)$ die Anzahl und $\sigma(n)$
die Summe der positiven Teiler von $n$.
Die beiden Größen werden hier nur für $n>0$ definiert.
\end{definition}

\begin{theorem}[Produktformeln für Teileranzahl und Teilersumme]
\label{b10guide:thm:divisor-product-formulas}
Sei $n=\prod_{i<r}p_i^{e_i}\in\mathbb N_+$ mit $r\in\mathbb N$,
paarweise verschiedenen Primzahlen $p_i$ und positiven natürlichen
Exponenten $e_i$. Dann gilt \cite{natarithfunctions}
\[
 \tau(n)=\prod_{i<r}(e_i+1),\qquad
 \sigma(n)=\prod_{i<r}\left(\sum_{j<e_i+1}p_i^j\right).
\]
Insbesondere gilt für jede Primzahl $p$ und jedes $e\in\mathbb N$
\[
 \tau(p^e)=e+1,\qquad
 \sigma(p^e)=\sum_{j<e+1}p^j.
\]
Für $n=1$ ist $r=0$; die leeren Produkte liefern
$\tau(1)=\sigma(1)=1$. Außerdem gilt für alle $n\in\mathbb N_+$
\[
 \tau(n)\text{ ist ungerade}
 \quad\Longleftrightarrow\quad
 \exists a\in\mathbb N:\ n=a^2.
\]
\end{theorem}

\begin{theorem}[Teileranzahl und Teilersumme eines teilerfremden Produkts]
\label{b10guide:thm:divisor-functions-coprime-product}
Für $a,b\in\mathbb N_+$ mit $\operatorname{ggT}(a,b)=1$ gilt
\cite{natarithfunctions}
\[
 \tau(ab)=\tau(a)\tau(b),\qquad
 \sigma(ab)=\sigma(a)\sigma(b).
\]
\end{theorem}

\section{Eulers \texorpdfstring{$\varphi$}{phi}-Funktion}

\begin{definition}[Eulers \texorpdfstring{$\varphi$}{phi}-Funktion]
\label{b10guide:def:totient}
Für $n\in\mathbb N_+$ und $k\in\mathbb N$ sei
\[
 \varepsilon_n(k)=
 \begin{cases}
  1,&\operatorname{ggT}(k,n)=1,\\
  0,&\operatorname{ggT}(k,n)\ne1.
 \end{cases}
\]
Eulers $\varphi$-Funktion ist für $n>0$ durch
\[
 \varphi(n)=\sum_{k<n}\varepsilon_n(k)
\]
festgelegt. Sie erfasst die Anzahl der zu $n$ teilerfremden
natürlichen Zahlen $k<n$, wobei $k=0$ zugelassen ist.
Insbesondere gilt $\varphi(1)=1$ \cite{nattotient}.
\end{definition}

\begin{theorem}[Produktformel für Eulers Funktion]
\label{b10guide:thm:totient-product-formula}
Sei $n=\prod_{i<r}p_i^{e_i}\in\mathbb N_+$ mit $r\in\mathbb N$,
paarweise verschiedenen Primzahlen $p_i$ und positiven natürlichen
Exponenten $e_i$. Dann gilt
\[
 \varphi(n)=\prod_{i<r}\bigl(p_i^{e_i-1}(p_i-1)\bigr).
\]
Für $n=1$ ist $r=0$ und der Produktwert ist $1$.
Für jede Primzahl $p$ und jedes $e\in\mathbb N_+$ gilt insbesondere
$\varphi(p^e)=p^{e-1}(p-1)$, also $\varphi(p)=p-1$
\cite{nattotient}.
\end{theorem}

\begin{theorem}[Eulers Funktion eines teilerfremden Produkts]
\label{b10guide:thm:totient-coprime-product}
Für $a,b\in\mathbb N_+$ mit $\operatorname{ggT}(a,b)=1$ gilt
\cite{nattotient}
\[
 \varphi(ab)=\varphi(a)\varphi(b).
\]
\end{theorem}

\begin{theorem}[Summe der Euler-Werte über die positiven Teiler]
\label{b10guide:thm:totient-divisor-sum}
Für jedes $n\in\mathbb N_+$ gilt mit
$\delta_n$ aus Definition~\ref{b10guide:def:divisor-functions} \cite{nattotient}
\[
 \sum_{i<n}\delta_n(i+1)\varphi(i+1)=n.
\]
\end{theorem}

\begin{theorem}[Eulers Satz über natürliche Restwerte]
\label{b10guide:thm:euler-totient}
Für $a\in\mathbb N$ und $m\in\mathbb N_+$ mit
$\operatorname{ggT}(a,m)=1$ gilt
\[
 a^{\varphi(m)}\mathbin{\mathrm{mod}}m
   =1\mathbin{\mathrm{mod}}m.
\]
Für $m\ge2$ ist der Rest also $1$; bei $m=1$ haben beide Seiten
den Wert $0$ \cite{natfermat}.
\end{theorem}


\section{Vollkommene natürliche Zahlen}

\begin{definition}[Vollkommene natürliche Zahl]
\label{b10guide:def:perfect-number}
Eine positive natürliche Zahl $n$ heißt \emph{vollkommen}, wenn
\[
 \sigma(n)=2n.
\]
Gleichwertig ist, daß die Summe ihrer positiven Teiler $d<n$
gleich $n$ ist. Die Teilersumme $\sigma$ ist in
Definition~\ref{b10guide:def:divisor-functions} festgelegt
\cite{natperfect}.
\end{definition}

\begin{theorem}[Euklid--Euler-Satz für gerade vollkommene Zahlen]
\label{b10guide:thm:euclid-euler-perfect}
Für $n\in\mathbb N_+$ gilt: Die Zahl $n$ ist genau dann gerade
und vollkommen, wenn es ein $k\in\mathbb N$ gibt, für das
\[
 2^{k+1}-1\text{ prim ist}
 \qquad\text{und}\qquad
 n=2^k(2^{k+1}-1).
\]
Die Primzahlbedingung erzwingt $k\ge1$. Die Aussage
charakterisiert die geraden vollkommenen Zahlen \cite{natperfect}.
\end{theorem}
